% Generated draft: requires mathematical and source review.
\documentclass[11pt]{article}
\usepackage[margin=25mm]{geometry}
\usepackage{fontspec}
\setmainfont{Latin Modern Roman}
\usepackage{amsmath,amssymb,mathtools,unicode-math}
\setmathfont{Latin Modern Math}
\usepackage{enumitem,needspace,xurl,hyperref}
\setlength{\parindent}{0pt}
\setlength{\parskip}{5pt}
\setlength{\emergencystretch}{3em}
\newcommand{\N}{\mathbb{N}}
\newcommand{\Z}{\mathbb{Z}}
\newcommand{\Q}{\mathbb{Q}}
\newcommand{\R}{\mathbb{R}}
\newcommand{\C}{\mathbb{C}}
\newcommand{\F}{\mathbb{F}}
\newcommand{\A}{\mathbb{A}}
\newcommand{\PP}{\mathbb{P}}
\newcommand{\E}{\mathbb{E}}
\newcommand{\eps}{\varepsilon}
\newcommand{\dd}{\,\mathrm{d}}
\begin{document}
\section*{UnsolvedMath ID: 11000205}
\textbf{Source problem code:} catalogue entry
\subsection{Definitions and mathematical statement}
For finitely generated groups $H\leq G$ with finite generating sets $B,A$, respectively, define
\[
 \Delta_H^G(n)=\max\{|h|_B:h\in H,\ |h|_A\leq n\}.
\]
The subgroup is undistorted if $|h|_B\leq C|h|_A+D$ for fixed constants $C,D$ and all $h\in H$; otherwise it is distorted. These notions are independent of finite generating-set choices up to the usual coarse comparison.

Problem 6 asks, for genus $g\geq3$, whether the Torelli subgroup of $M_{g,m}$ is undistorted, and more generally asks for a distorted finitely generated subgroup of $M_{g,m}$. In the customary convention the Torelli group is the kernel of the action on integral genus-$g$ surface homology. The meaning of $m$, whether homology is taken before filling marked points, and the exact closed/marked surface conventions are absent.

The registry marks the question solved negatively and attributes to Broaddus--Farb--Putman an exponential lower bound and a doubly exponential upper bound for the distortion of $\mathcal I_g$ and $\mathcal I_{g,1}$ for $g\geq3$. These stated bounds imply distortion but do not by themselves give matching exponential upper and lower bounds, despite the summary's shorter claim ``exponentially distorted.'' Their extension to every interpretation of $m$ is not supplied.

The following algorithmic material is context. The source describes train-track splitting sequences with uniformly bounded choices at each step as a means of treating mapping-class uniform quasigeodesics algorithmically. It recalls the word problem (deciding whether an input word is the identity) and conjugacy problem (deciding whether $h=ugu^{-1}$ for some $u$). It then describes a biautomatic structure: a finite alphabet $A$ with symmetric generating map $\pi:A\to\Gamma$, an involution $\iota$ satisfying $\pi(\iota a)=\pi(a)^{-1}$, and a regular language $L\subset A^*$ mapping onto $\Gamma$, with a uniform synchronous fellow-traveller bound for paths represented by $xwy$ and normal words $w'\in L$ representing the same element, for $x,y\in A$, $w\in L$. Since path lengths differ, endpoint-extension and time conventions must be specified; they are omitted in the extraction. No explicit complexity bound is given. The passage attributes this framework to Epstein, Cannon, Holt, Levy, Paterson and Thurston [E], mentions Mosher [M1] and [H1], and ends before the announced theorem. No biautomaticity theorem is asserted here.
\subsection{Short English statement}
Is the Torelli subgroup linearly embedded in the mapping class group, and can one find distorted finitely generated subgroups? The registry reports a negative answer with exponential lower and doubly exponential upper distortion bounds in specified Torelli cases.
\subsection{Sources}
\begin{itemize}
\item Local registry: \url{data/unsolvedmath/problems.json}, permanent UnsolvedMath ID $11000205$, exact source code \texttt{catalogue entry}.
\item Farb, \emph{Problems on Mapping Class Groups and Related Topics}, Problem 6, PDF page 249: \url{https://www.math.uchicago.edu/~farb/papers/mcgbook.pdf}. Supplied access date: 2026-07-29; extraction: pdf-font-structure; repeated local numbering disambiguated by page among 329 headings.
\item Supplied solution reference: Broaddus--Farb--Putman, ``Irreducible Sp-representations and subgroup distortion in the mapping class group,'' \emph{Comm. Math. Helv.} 86 (2011), associated with \url{https://arxiv.org/abs/1006.0787}. The summary attributes the question to Hamenst\"adt and earlier rank-one/recursive-distortion context to Farb--Lubotzky--Minsky. Bibliographic matching and precise hypotheses require verification.
\item Supplied contextual references: [BH], Chapter III.$\Gamma$.1; Epstein--Cannon--Holt--Levy--Paterson--Thurston [E]; Mosher [M1]; [H1]. Full bibliography entries and the announced following theorem are absent.
\end{itemize}
\end{document}
